\section*{Multivibratore Triangolare}

\noindent\colorbox{orange}{
  \begin{minipage}{1.0\linewidth}
    \subsection*{Richiesta (blocco 1, a)}
    Dalle caratteristiche dei segnali osservazioni ai punti 1a, 1d, 1e, descrivere la funzione
    svolta della prima parte del circuito (blocco 1 in figura, ingresso $V_1$, uscita $V_{TR}$).
  \end{minipage}
}
\subsection*{Risposta (blocco 1, a)}

\noindent\colorbox{orange}{
  \begin{minipage}{1.0\linewidth}
    \subsection*{Richiesta (blocco 1, b)}
    Determinare analiticamente la funzione di trasferimento ed individuarne le
    caratteristiche (poli, zeri) nell'ipotesi di idealit\'a di $U_1$.
  \end{minipage}
}
\subsection*{Risposta (blocco 1, b)}
L'ingresso \textit{invertente} ideale avendo impedenza ``grande'' \'e soggetto alla tensione
del \textit{partitore di tensione} $R_1R_2$
\begin{equation*}
  V_+ = V_{TR}/2
\end{equation*}
L'ingresso \textit{\underline{non}-invertente} ideale avendo impedenza ``grande'' \'e soggetto alla
tensione compatibile con la KCL 
\begin{equation*}
  \begin{gathered}
    \frac{V_1 - V_-}{Z_{R_3}} = \frac{V_-}{Z_{C_1}} + \frac{V_- - V_{TR}}{Z_{R_4}} \\ \implies
    V_- = \frac{V_1 + V_{TR}}{2 + sRC}
  \end{gathered}
\end{equation*}
sott'ipotesi di op-amp ideale (con guadagno differenziale ``grande'' e regime di operazione lineare)
applico il \textit{principio di corto-circuito virtuale}
\begin{equation*}
  \begin{gathered}
    V_+ = V_- \\ \implies
    \frac{V_{TR}}{2} = \frac{V_1 + V_{TR}}{2 + sRC}
  \end{gathered}
\end{equation*}
la funzione di trasferimento complessa si ottiene dal rapporto $V_{TR}/V_1$
\begin{equation*}
  \colorbox{blu}{\boxed{H_1(s) = \frac{2}{sRC}}}
\end{equation*}
non ci sono \textit{zeri} ma c\'e un unico \textit{polo} per $s = 0$. \\

\noindent\colorbox{orange}{
  \begin{minipage}{1.0\linewidth}
    \subsection*{Richiesta (blocco 1, c)}
    Sulla base di quanto ottenuto al punto precedente, discutere la risposta dello stesso
    blocco nel dominio del tempo; di che tipo di circuito si tratta?
  \end{minipage}
}
\subsection*{Risposta (blocco 1, c)}
La funzione di trasferimento complessa per il circuito \'e quella di un'integratore con
\textit{centro-banda} 0 e \textit{guadagno massimo} 2, nel dominio del tempo per mezzo
della \textit{trasformata di laplace inversa} si ottiene
\begin{equation*}
  \colorbox{blu}{\boxed{V_{TR}(t) = \frac{2}{RC}\int_0^t V_1(t')~\text{d}t'}}
\end{equation*}

\noindent\colorbox{orange}{
  \begin{minipage}{1.0\linewidth}
    \subsection*{Richiesta (blocco 1, d)}
    Dedurre una relazione quantitativa tra i livelli del segnale di ingresso ($V_{SQ}$) e la
    pendenza dell'uscita ($V_{TR}$), confrontandola con quanto osservato/misurato ai precedenti
    punti 1b, 1c.
  \end{minipage}
}
\subsection*{Risposta (blocco 1, d)}
Per $V_{SQ}$ segnale di un'onda quadra la tensione $V_1 \equiv \text{const.}$ e pu\'o essere portata
fuori dall'integrale diventando cos\'i una funzione lineare la cui pendenza
\begin{equation*}
  \colorbox{blu}{\boxed{\frac{\text{d}}{\text{d}t}V_{TR} = \frac{2V_{SQ}}{RC}}}
\end{equation*}


\noindent\colorbox{orange}{
  \begin{minipage}{1.0\linewidth}
    \subsection*{Richiesta (blocco 2, a)}
    Descrivere qualitativamente la rete identificata in figura come blocco 2, eventualmente
    riconducendola ad uno schema gi\'a noto.
  \end{minipage}
}
\subsection*{Risposta (blocco 2, a)}
Il secondo blocco costituisce un \textit{trigger di Schmitt}. \\

\noindent\colorbox{orange}{
  \begin{minipage}{1.0\linewidth}
    \subsection*{Richiesta (blocco 2, b)}
    Verificato che $V_{SQ}$ si attesti sui livelli di saturazione di $U_2$, determinarne le soglie,
    ovvero i valori del suo ingresso $V_{TR}$ in corrispondenza delle quali si verifica la
    commutazione del suo stato d'uscita e confrontarli con l'ampiezza delle rampe di $V_{TR}$
    misurata al punto 1c.
  \end{minipage}
}
\subsection*{Risposta (blocco 2, b)}
\begin{equation*}
  \begin{gathered}
    V_{thr-} = V_{OL}/2 \\
    V_{thr+} = V_{OH}/2
  \end{gathered}
\end{equation*}

\noindent\colorbox{orange}{
  \begin{minipage}{1.0\linewidth}
    \subsection*{Richiesta (blocco 2, c)}
    Mostrare che, per effetto della connessione ad anello dei due blocchi, entrambi gli stati
    di uscita di $V_{SQ}$ sono metastabili.
  \end{minipage}
}
\subsection*{Risposta (blocco 2, c)}
\begin{itemize}
\item
  Il segnale entrante nel primo blocco \'e un segnale costante, partendo da $t = 0$ esso
  crescer\'a/decrescer\'a in base al segno del segnale d'entrata.
\item
  Il segnale crescer\'a/decrescer\'a fino a superare la soglia del secondo blocco, in
  quell'istante il segnale in uscita dal secondo blocco e quindi entrante al primo cambier\'a
  di segno.
\end{itemize}
L'azione di \textit{inversione} del segnale causata dal secondo blocco assieme alla \textit{crescita costante}
(in modulo) del primo blocco fa s\'i che lo stato di entrambi i segnali commuti periodicamente. \\

\noindent\colorbox{orange}{
  \begin{minipage}{1.0\linewidth}
    \subsection*{Richiesta (blocco 2, d)}
    Nell'ipotesi di idealità di $U_2$ e di simmetria dei suoi livelli di saturazione, calcolare il
    periodo e confrontarlo con quanto misurato al punto 1c.
  \end{minipage}
}
\subsection*{Risposta (blocco 2, d)}
\begin{equation*}
  \frac{V_{OH}}{2} = \frac{2V_{OH}}{RC}t - \frac{V_{OH}}{2} \implies t = RC/2
\end{equation*}
\begin{equation*}
  \frac{V_{OL}}{2} = \frac{2V_{OL}}{RC}t - \frac{V_{OL}}{2}
\end{equation*}
se i livelli di saturazione sono simmetrici
\begin{equation*}
  -\frac{V_{OH}}{2} = -\frac{2V_{OH}}{RC}t + \frac{V_{OH}}{2} \implies t = RC/2
\end{equation*}
il periodo \'e dunque
\begin{equation*}
  \colorbox{blu}{\boxed{T = RC}}
\end{equation*}
